\subsection*{Tvrdý specializační dril}
\addcontentsline{toc}{subsection}{Tvrdý specializační dril}
\begin{infobox}[title={Pokyny}]
\footnotesize 6 specializačních otázek za sebou (numerické). Vázající podlaha je specializace ($\ge7/12$, tři ze čtyř na $\ge2$). Trénuj rychlost a přesnost výpočtu.
\end{infobox}

\mockq{Otázka 1}{specializace}{EM: responsibility}
\begin{zadani}$\pi=(0{,}5,0{,}5)$, $p(x{=}1\mid1)=0{,}7$, $p(x{=}1\mid2)=0{,}3$. Spočtěte responsibility bodu $x{=}1$.\end{zadani}
\begin{reseni}$r_1=\dfrac{0{,}5\cdot0{,}7}{0{,}5\cdot0{,}7+0{,}5\cdot0{,}3}=\dfrac{0{,}35}{0{,}5}=0{,}7$, $r_2=0{,}3$.\end{reseni}

\mockq{Otázka 2}{specializace}{Matice záměn}
\begin{zadani}$TP{=}30,FP{=}10,FN{=}20,TN{=}40$. Spočtěte precision, recall, $F_1$ a accuracy.\end{zadani}
\begin{reseni}precision $=\tfrac{30}{40}=0{,}75$, recall $=\tfrac{30}{50}=0{,}6$, $F_1=\tfrac{2\cdot0{,}75\cdot0{,}6}{1{,}35}\approx0{,}667$, accuracy $=\tfrac{70}{100}=0{,}7$.\end{reseni}

\mockq{Otázka 3}{specializace}{HMM: filtrace}
\begin{zadani}Prior $(0{,}6,0{,}4)$; přechod $P(t\mid t)=0{,}7,P(t\mid f)=0{,}2$; emise $P(o\mid t)=0{,}8,P(o\mid f)=0{,}1$. Spočtěte $P(\cdot\mid o)$.\end{zadani}
\begin{reseni}Predikce: $P(t)=0{,}7\cdot0{,}6+0{,}2\cdot0{,}4=0{,}5$, $P(f)=0{,}5$. Korekce: $\propto(0{,}8\cdot0{,}5,\,0{,}1\cdot0{,}5)=(0{,}4,\,0{,}05)$, norm $/0{,}45\to(0{,}889,\,0{,}111)$.\end{reseni}

\mockq{Otázka 4}{specializace}{Bayes-optimal predikce}
\begin{zadani}$h_1{:}P(+)=0{,}9$ (prior $0{,}4$), $h_2{:}P(+)=0{,}5$ (prior $0{,}6$). Po pozorování ,,$+$'' spočtěte posterior a $P(\text{další}{=}+)$.\end{zadani}
\begin{reseni}Posterior $\propto(0{,}4\cdot0{,}9,\,0{,}6\cdot0{,}5)=(0{,}36,\,0{,}30)$, norm $/0{,}66\to(0{,}545,\,0{,}455)$. Predikce $=0{,}9\cdot0{,}545+0{,}5\cdot0{,}455\approx0{,}718$.\end{reseni}

\mockq{Otázka 5}{specializace}{Ensemble: majoritní hlasování}
\begin{zadani}Tři nezávislé klasifikátory, každý přesnost $0{,}7$, majoritní hlasování. Spočtěte přesnost a uveďte nejlepší/nejhorší případ.\end{zadani}
\begin{reseni}Při nezávislosti $P(\ge2\text{ správně})=\binom32 0{,}7^2\cdot0{,}3+0{,}7^3=3\cdot0{,}147+0{,}343=0{,}784$. Nejlepší případ (chyby se nepřekrývají) až $1$; jsou-li chyby dokonale korelované, ensemble nepomůže a zustává $0{,}7$; při nejnepříznivější závislosti (maximální překryv dvojic chyb) může klesnout až k $0{,}55$. Ensemble pomáhá při \emph{nekorelovaných} chybách.\end{reseni}

\mockq{Otázka 6}{specializace}{Iterace hodnot}
\begin{zadani}Stavy $A,B$, $\gamma=0{,}9$. Z $A$ akce do $B$ (odměna 0), z $B$ akce do cíle (odměna 10). Spočtěte $V^*(A),V^*(B)$.\end{zadani}
\begin{reseni}$V^*(B)=10+0{,}9\cdot0=10$. $V^*(A)=0+0{,}9\cdot V^*(B)=0{,}9\cdot10=9$. (Iterace: $V_1(B)=10$, $V_2(A)=9$.)\end{reseni}

\begin{infobox}[title={Vyhodnocení}]
\footnotesize Ber to jako dva specializační bloky po čtyřech. V každém čtyřbloku miř na $\ge7$: tři otázky vyřešené úplně ($3{\times}2$) plus jedna alespoň částečně ($1$) stačí. Pokud ti numerika padá pod 2 body, zopakuj příslušnou úlohu z úrovní 1-2.
\end{infobox}
